\section{算法设计}

\subsection{共同基础：进化操作设计}

三种并行算法共享相同的进化操作框架，本节对所使用的编码方案、适应度函数及遗传算子进行统一说明。

\subsubsection{编码与适应度函数}

采用最直观的路径表示法（Permutation Encoding）对 TSP 解进行编码。每个个体由一个长度为 442 的全排列 $\pi = (\pi_1, \pi_2, \dots, \pi_{442})$ 表示，其中 $\pi_i$ 为城市编号，排列的顺序即为旅行商访问城市的顺序。路径的总长度为

\[
L(\pi) = \sum_{i=1}^{441} d_{\pi_i,\pi_{i+1}} + d_{\pi_{442},\pi_1},
\]

其中 $d_{ij}$ 为城市 $i$ 与城市 $j$ 之间的欧氏距离。适应度函数定义为路径长度的倒数：

\[
f(\pi) = \frac{1}{L(\pi)},
\]

路径越短则适应度越高，选择操作倾向于保留适应度较高的个体。

\subsubsection{种群初始化}

初始种群由随机生成的全排列构成。对于每个个体，采用 Fisher--Yates 洗牌算法生成长度为 442 的随机排列，确保搜索空间初始覆盖的均匀性。dEA 的全局种群规模为 200，HdEA 与 HdEA-MP 的每个岛屿种群规模为 50。

\subsubsection{选择算子}

采用锦标赛选择（Tournament Selection）作为选择算子。每代在选择父代时，从当前种群中随机抽取 $k=5$ 个个体，比较其适应度，选择适应度最高的个体进入交配池。锦标赛选择的优点是计算简单、不依赖全局适应度排序、且可以通过调整锦标赛规模 $k$ 灵活控制选择压力——$k$ 越大，选择压力越强，种群收敛越快但也更容易陷入局部最优。

\subsubsection{交叉算子}

采用顺序交叉（Order Crossover, OX）作为交叉算子。给定两个父代个体 $P_1$ 和 $P_2$，OX 的操作步骤如下：

\begin{enumerate}
    \item 在排列长度范围内随机选择两个切点 $a$ 和 $b$（$a < b$），将父代划分为左、中、右三段。
    \item 子代 $O_1$ 继承父代 $P_1$ 的中间片段 $P_1[a:b]$，保持位置不变。
    \item 从父代 $P_2$ 的 $b$ 位置开始循环遍历，将所有不在 $O_1$ 中的城市按顺序填入 $O_1$ 的剩余空位（从位置 $b$ 开始，循环至位置末尾再回到开头）。
    \item 对称地，子代 $O_2$ 以 $P_2$ 的中间片段为基础，从 $P_1$ 中填充剩余城市。
\end{enumerate}

OX 的独特优势在于它不仅继承了两个父代的基因片段，还保留了父代中城市之间的相对顺序信息，同时保证子代仍然是合法的全排列。对于 442 规模的 TSP，OX 能够在路径重构与节段继承之间取得良好平衡。

\subsubsection{变异算子}

采用交换变异（Swap Mutation）。对每个新生成的子代，以概率 $p_m = 0.02$ 随机选择其编码中的两个不同位置，交换这两个位置上的城市编号。变异操作的主要作用是在进化后期种群多样性下降时，为种群引入微小的扰动，帮助算法跳出局部最优。

\subsubsection{精英保留策略}

每代在完成选择、交叉和变异后，从上一代种群中选取适应度最高的 3 个个体（精英），直接复制到下一代种群中，替换掉子代种群中最差的 3 个个体。精英保留策略保证了算法的最优解在世代间不会退化，是进化算法收敛性的重要保障。

\subsection{Algorithm 1：主从式 dEA}

主从式分布式进化算法（dEA）采用最简单的并行架构——一个 Master 进程负责全局进化逻辑，多个 Worker 进程负责最耗时的适应度评估环节。这种架构的核心思想是将计算密集部分并行化，而保持控制逻辑的集中以提高协调效率。

\subsubsection{架构与通信模式}

dEA 使用 MPI\_COMM\_WORLD 中的全部 4 个进程，其中 rank 0 为 Master，rank 1--3 为 Worker。Master 维护一个全局种群（规模 POP = 200），在每一代中依次执行以下步骤：

\begin{enumerate}
    \item \textbf{选择与繁殖：}Master 执行锦标赛选择、OX 交叉和交换变异，从上一代的 200 个个体中生成 200 个全新子代。
    \item \textbf{分发评估：}Master 利用 MPI\_Bcast 将 200 个子代组成的数组广播至所有 Worker。每个 Worker 分到约 67 个子代个体。
    \item \textbf{并行评估：}3 个 Worker 各自计算所分配子代的 TSP 路径距离（即适应度），评估过程完全独立、无进程间通信。
    \item \textbf{结果收集：}Worker 通过 MPI\_Gatherv 将评估结果（距离值数组）发送回 Master。Master 根据收集到的适应度值更新全局种群。
\end{enumerate}

此过程重复执行直到达到最大进化代数 MAX\_GEN = 200000。

\subsubsection{算法伪码}

\begin{lstlisting}[language={},basicstyle=\scriptsize\fontspec{DejaVu Sans Mono}]
Algorithm 1: dEA (Master-Worker)
  if rank == 0 then           // Master
      POP ← 200, MAX_GEN ← 200000
      初始化全局种群 pop[POP]
      for gen = 1 to MAX_GEN do
          offspring ← SelectionCrossoverMutation(pop)
          MPI_Bcast(offspring, 0, MPI_COMM_WORLD)
          MPI_Gatherv(fitness, 0, MPI_COMM_WORLD)
          pop ← UpdatePopulation(pop, offspring, fitness)
      end for
  else                         // Worker
      while true do
          MPI_Bcast(offspring, 0, MPI_COMM_WORLD)
          for each ind in offspring do
              fitness[ind] ← EvaluateTSPDistance(ind)
          end for
          MPI_Gatherv(fitness, 0, MPI_COMM_WORLD)
      end while
  end if
\end{lstlisting}

\subsubsection{通信复杂度分析}

dEA 每代需要两次全局集体通信操作：MPI\_Bcast 广播子代数组（$\mathcal{O}(\text{POP} \cdot n)$ 字节）和 MPI\_Gatherv 收集评估结果（$\mathcal{O}(\text{POP})$ 字节），其中 $n=442$ 为城市数。在 $n_{\text{np}} = 4$ 的条件下，Bcast 的通信复杂度为 $\mathcal{O}(\text{POP} \cdot n \cdot \log n_{\text{np}} / n_{\text{np}})$。考虑到 Master 上的选择交叉操作与 Worker 上的适应度评估之间存在天然流水线重叠，实际运行中通信开销被有效隐藏。计算负载方面，由于子代被均匀分发，各 Worker 承担的评估量相等（偏差 $\le 1$），实现了高度均衡的负载分配。

\subsection{Algorithm 2：分层 dEA（HdEA）}

HdEA 将种群划分为多个独立演化的岛屿，并在岛屿之上增加一层精英协调，形成两层级联式架构。其设计动机源于岛屿模型在维持种群多样性方面的理论优势——通过限制个体交配的范围，不同岛屿可以探索搜索空间中差异化的区域，延缓早熟收敛。

\subsubsection{架构与通信模式}

HdEA 同样使用 4 个 MPI 进程，但每个进程的角色与 dEA 有本质区别：

\begin{itemize}
    \item \textbf{第一层（岛屿并行层）：}4 个进程各自维护一个完全独立的子种群（规模 POP\_ISLAND = 50），在各进程内部独立执行完整的 EA 流程——锦标赛选择、OX 交叉、交换变异、适应度评估和精英保留——持续运行 MIG\_FREQ = 20 代。在此阶段，进程之间完全不存在任何通信，每个岛屿相当于一个串行 EA。

    \item \textbf{第二层（精英协调层）：}每隔 MIG\_FREQ 代，各进程通过 MPI\_Gather 将本岛适应度最高的 3 个精英个体发送至 Master（rank 0）。Master 将所有岛屿的精英汇总（共 $4 \times 3 = 12$ 个个体），选择其中全局最优的若干个体，再通过 MPI\_Bcast 广播给所有进程。各进程据此更新本地的精英存档。
\end{itemize}

\subsubsection{设计动机与预期效果}

HdEA 的设计基于以下关键洞察：种群多样性对于进化算法的搜索能力至关重要。根据遗传漂变理论，小种群的遗传多样性会因采样效应快速丢失，而将一个大种群隔离为多个小型子种群，可以通过限制基因流动来减缓局部最优的传播。理想情况下，不同岛屿在各自的小生态位中独立进化，可能发现不同的高质量解区域，再通过周期性的精英迁移实现全局信息的共享。

然而，这种设计方案也面临一个内在矛盾：岛屿规模越小，多样性损失越快，但岛屿数量越多，并行度越高。在 4 进程的条件下，每个岛屿只能分到 50 个个体，对于一个 442 城市的搜索空间而言，50 个个体所能覆盖的搜索范围极为有限。

\subsection{Algorithm 3：带种群迁移的 HdEA（HdEA-MP）}

HdEA-MP 在 HdEA 的两层架构之上，额外引入第三层机制——环形种群迁移（Ring Migration），通过增加岛屿之间的个体交换频率和交换量，缓解岛屿规模过小导致的多样性不足问题。

\subsubsection{环形迁移机制}

与 HdEA 仅在 Master 处进行精英汇聚不同，HdEA-MP 在每 MIG\_FREQ = 20 代时，除执行精英迁移外，额外执行一次进程间的直接个体交换。具体流程如下：

\begin{enumerate}
    \item 每个进程从自身子种群中挑选适应度最高的 MIG\_M = 3 个个体。
    \item 将这 3 个个体的完整路径编码（$3 \times 442$ 个整数）和适应度值（3 个双精度浮点数，但发送时取整为整数）打包为一个平面缓冲区。
    \item 通过 MPI\_Sendrecv 将打包数据发送至右邻居进程 $(rk + 1) \bmod 4$，同时从左邻居进程 $(rk - 1 + 4) \bmod 4$ 接收同样格式的迁移数据。
    \item 对接收到的 3 个迁移个体，逐一判断其路径长度是否小于本地种群中的最差个体（即适应度是否更高），若是，则用迁移个体替换最差个体。
\end{enumerate}

环形拓扑在并行计算中具有天然的对称性和良好的扩展性。对于 $n_{\text{np}}$ 个进程，环形通信仅需 $n_{\text{np}}$ 个点到点通信（或 $n_{\text{np}}$ 对 MPI\_Sendrecv），通信复杂度为 $\mathcal{O}(n_{\text{np}})$，且不需要 Master 中转，避免了集中式通信的瓶颈。

\subsubsection{MPI\_Sendrecv 的设计选择}

在最初的原型实现中，环形迁移采用交替的 MPI\_Send 和 MPI\_Recv 进行通信。但测试中发现，当 MIG\_M = 3 时，发送方和接收方的 MPI 内部缓冲区可能同时达到饱和，导致死锁——发送操作阻塞等待接收缓冲区释放，而接收操作阻塞等待发送完成。解决这一问题有两种途径：一是使用 MPI\_Sendrecv 替代单独的 Send 和 Recv，二是使用非阻塞通信 MPI\_Isend / MPI\_Irecv。本实现选择 MPI\_Sendrecv，因为它以单一函数调用的方式同时处理发送和接收，由 MPI 实现内部协调缓冲区的使用，代码简洁且不易出错。

\subsubsection{算法伪码}

\begin{lstlisting}[language={},basicstyle=\scriptsize\fontspec{DejaVu Sans Mono}]
Algorithm 3: HdEA-MP (Ring Migration)
  POP_ISLAND ← 50, MAX_GEN ← 200000
  MIG_FREQ ← 20, MIG_M ← 3
  初始化本地子种群 pop[POP_ISLAND]
  for gen = 1 to MAX_GEN do
      // 本地进化
      pop ← SelectionCrossoverMutation(pop)
      
      if gen mod MIG_FREQ == 0 then
          // 第1层：精英迁移（同 HdEA）
          MPI_Gather(local_best, 0, MPI_COMM_WORLD)
          if rank == 0 then 更新全局精英池 end if
          MPI_Bcast(global_elite, 0, MPI_COMM_WORLD)
          
          // 第2层：环形迁移
          sendTo ← (rank + 1) mod 4
          recvFrom ← (rank - 1 + 4) mod 4
          打包 MIG_M 个最优个体至 sendBuf
          MPI_Sendrecv(sendBuf, sendTo, recvBuf, recvFrom)
          用 recvBuf 中的更优个体替换本地最差个体
      end if
  end for
\end{lstlisting}

\subsubsection{与 HdEA 的本质区别}

HdEA-MP 比 HdEA 多出的一层环形迁移带来两个重要变化：（1）迁移量从 HdEA 的每个岛屿仅贡献一个精英个体变为每个岛屿交换三个个体，基因交流的密度大幅提高；（2）迁移发生在岛屿之间（直接点对点），而非通过 Master 中转，因此环形迁移可以在不增加 Master 负担的情况下进行。然而，这种密集迁移也是一把双刃剑——如果迁移频率过高或迁移量过大，各岛屿的种群将被强制拉近，丧失探索差异区域的能力。
